\documentclass[12pt,a4paper]{article}

\usepackage[utf8]{inputenc}
\usepackage[T1]{fontenc}
\usepackage{lmodern}
\usepackage{textcomp}
\usepackage{newunicodechar}
\usepackage[margin=2.5cm]{geometry}
\usepackage{amsmath,amssymb}
\usepackage{graphicx}
\usepackage{listings}
\usepackage{xcolor}
\usepackage{hyperref}
\usepackage{fancyhdr}
\usepackage{float}
\usepackage{booktabs}
\usepackage{caption}
\usepackage{tabularx}
\usepackage{array}

% Adjust header height for fancyhdr
\setlength{\headheight}{15pt}
\addtolength{\topmargin}{-3pt}

% Romanian Unicode support for standard pdflatex
\newunicodechar{Ă}{\v{A}}
\newunicodechar{ă}{\v{a}}
\newunicodechar{Â}{\^{A}}
\newunicodechar{â}{\^{a}}
\newunicodechar{Î}{\^{I}}
\newunicodechar{î}{\^{i}}
\newunicodechar{Ș}{\c{S}}
\newunicodechar{ș}{\c{s}}
\newunicodechar{Ț}{\c{T}}
\newunicodechar{ț}{\c{t}}

\hypersetup{
    colorlinks=true,
    linkcolor=blue,
    filecolor=magenta,
    urlcolor=cyan,
    citecolor=blue
}

% Colors for code listings
\definecolor{codebg}{gray}{0.96}
\definecolor{codegreen}{rgb}{0.13,0.55,0.13}
\definecolor{codegray}{rgb}{0.5,0.5,0.5}
\definecolor{codepurple}{rgb}{0.58,0,0.82}
\definecolor{codeblue}{rgb}{0.05,0.25,0.70}

\lstset{
    backgroundcolor=\color{codebg},
    basicstyle=\ttfamily\footnotesize,
    breaklines=true,
    breakatwhitespace=true,
    frame=single,
    rulecolor=\color{gray!40},
    keepspaces=true,
    showstringspaces=false,
    upquote=true,
    columns=fullflexible,
    tabsize=2
}

\lstdefinestyle{cstyle}{
    language=C,
    commentstyle=\color{codegreen}\itshape,
    keywordstyle=\color{codeblue}\bfseries,
    stringstyle=\color{codepurple},
    numbers=left,
    numbersep=7pt,
    numberstyle=\tiny\color{codegray},
    stepnumber=1
}

\lstdefinestyle{bashstyle}{
    language=bash,
    commentstyle=\color{codegreen}\itshape,
    keywordstyle=\color{codeblue}\bfseries,
    stringstyle=\color{codepurple},
    numbers=left,
    numbersep=7pt,
    numberstyle=\tiny\color{codegray},
    stepnumber=1
}

\lstdefinestyle{consolestyle}{
    backgroundcolor=\color{codebg},
    basicstyle=\ttfamily\footnotesize,
    breaklines=true,
    frame=single,
    rulecolor=\color{gray!40},
    numbers=none
}

% Header and Footer
\pagestyle{fancy}
\fancyhf{}
\lhead{Operating Systems -- Lab 9}
\rhead{Iulian Pavlov, FAF-241}
\cfoot{\thepage}

\begin{document}

% ============================================================================
% TITLE PAGE
% ============================================================================
\hypersetup{pageanchor=false}
\begin{titlepage}
\begin{center}
    \includegraphics[width=0.3\textwidth]{logo.png} \\[0.2cm]

    \textbf{MINISTRY OF EDUCATION, CULTURE AND RESEARCH \\ OF THE REPUBLIC OF MOLDOVA} \\[0.3cm]

    \textbf{Technical University of Moldova \\ Faculty of Computers, Informatics and Microelectronics \\ Department of Software and Automation Engineering} \\[1.8cm]

    \textbf{Iulian Pavlov, FAF-241}\\[0.5cm]

    \Huge \textbf{Report} \\[0.5cm]

    \large Laboratory work nr. 9 \\[0.5cm]

    \textbf{\Large Watching System Calls with \texttt{strace}} \\[0.3cm]
    \normalsize Course: Operating Systems (Track A: the real OS, from the outside (Linux)) \\[2.2cm]

    \begin{flushright}
        \textit{Checked by:} \\
        \textbf{Reițman Patricia}, \textit{University Assistant} \\
        Department of Software and Automation Engineering, \\
        FCIM Faculty, UTM
    \end{flushright}
    \vfill

    Chișinău -- 2026
\end{center}
\end{titlepage}
\hypersetup{pageanchor=true}

\pagenumbering{roman}
\newpage
\tableofcontents
\newpage
\pagenumbering{arabic}

% ============================================================================
% 1. THEORETICAL BACKGROUND & OBJECTIVES
% ============================================================================
\section{Theoretical Background and Objectives}

\subsection{Laboratory Objectives and Scope}
Every interaction between an unprivileged user-space application and the Linux operating system kernel passes through the \textbf{system call interface}. Ordinary library routines execute entirely within the process's virtual address space, whereas system calls cross the hardware privilege boundary into kernel mode to perform file I/O, directory traversal, memory allocation, and process management.

The objectives of Laboratory Work Nr.~9 are:
\begin{itemize}
    \item \textbf{Inspecting Real Commands with \texttt{strace}}: Tracing the file and I/O system calls (\texttt{openat}, \texttt{read}, \texttt{write}, \texttt{close}) executed by \texttt{ls -l}, and profiling its complete system call distribution using \texttt{strace -c}.
    \item \textbf{Direct Kernel I/O in C (\texttt{syscalls.c})}: Writing, compiling, and executing a C program that performs file creation, writing, reading, and closing using raw POSIX system calls (\texttt{open}, \texttt{write}, \texttt{read}, \texttt{close}) instead of buffered \texttt{<stdio.h>} file streams.
    \item \textbf{Understanding File Descriptor Allocation}: Explaining why the first user-opened file receives descriptor \texttt{3} rather than \texttt{0}, \texttt{1}, or \texttt{2}.
    \item \textbf{Quantifying System Calls and Buffering}: Counting specific calls (\texttt{write}, \texttt{openat}) across commands and examining how \texttt{libc} stream buffering reduces kernel crossings.
\end{itemize}

\subsection{System Call Mechanics and File Descriptors}
\begin{enumerate}
    \item \textbf{Tracing via \texttt{ptrace}}: The \texttt{strace} utility uses the Linux \texttt{ptrace()} system call to attach to a child process, intercepting each system call entry and exit to decode registers, arguments, and return codes (including negative \texttt{errno} errors such as \texttt{ENOENT}).
    \item \textbf{Standard File Descriptors}: Every newly created POSIX process begins with three pre-opened descriptors inherited from its parent shell: \texttt{0} (\texttt{STDIN\_FILENO}, standard input), \texttt{1} (\texttt{STDOUT\_FILENO}, standard output), and \texttt{2} (\texttt{STDERR\_FILENO}, standard error). The kernel always assigns the lowest unused integer descriptor when opening a new file.
    \item \textbf{C Wrapper \texttt{open()} vs.\ Kernel \texttt{openat()}}: On modern Linux systems (such as x86\_64 with glibc), the C function \texttt{open(path, flags, mode)} is implemented in user space as a thin wrapper around the \texttt{openat(AT\_FDCWD, path, flags, mode)} kernel system call, where \texttt{AT\_FDCWD} instructs the kernel to resolve relative paths against the current working directory.
    \item \textbf{Direct Syscalls vs.\ \texttt{libc} Buffering}: Calling \texttt{write()} transfers bytes directly to the kernel on every invocation. By contrast, \texttt{printf()} writes into a user-space \texttt{stdout} buffer: connected to a terminal (TTY), \texttt{stdout} is \emph{line-buffered} (flushing on each \texttt{\textbackslash n}); redirected to a pipe, file, or \texttt{/dev/null}, it switches to \emph{block buffering} and flushes multiple lines together.
\end{enumerate}

\newpage

% ============================================================================
% 2. PART 1: SEE THE CALLS OF A REAL COMMAND
% ============================================================================
\section{Part 1: Watching System Calls of a Real Command (\texttt{ls -l})}

\subsection{Tracing File and I/O Operations}
We traced the open, read, write, and close system calls issued during a single execution of \texttt{ls -l}:
\begin{lstlisting}[style=consolestyle]
$ strace -f -e trace=openat,read,write,close ls -l
openat(AT_FDCWD, "/etc/ld.so.cache", O_RDONLY|O_CLOEXEC) = 3
close(3)                                = 0
openat(AT_FDCWD, "/usr/lib/libcap.so.2", O_RDONLY|O_CLOEXEC) = 3
read(3, "\177ELF\2\1\1\0\0\0\0\0\0\0\0\0\3\0>\0\1\0\0\0\320{\0\0\0\0\0\0"..., 1024) = 1024
close(3)                                = 0
openat(AT_FDCWD, "/usr/lib/libc.so.6", O_RDONLY|O_CLOEXEC) = 3
read(3, "\177ELF\2\1\1\3\0\0\0\0\0\0\0\0\3\0>\0\1\0\0\0\300y\2\0\0\0\0\0"..., 1024) = 1024
close(3)                                = 0
...
openat(AT_FDCWD, ".", O_RDONLY|O_NONBLOCK|O_CLOEXEC|O_DIRECTORY) = 3
openat(AT_FDCWD, "/etc/nsswitch.conf", O_RDONLY|O_CLOEXEC) = 4
read(4, "# Name Service Switch configurat"..., 4096) = 372
read(4, "", 4096)                       = 0
close(4)                                = 0
openat(AT_FDCWD, "/etc/passwd", O_RDONLY|O_CLOEXEC) = 4
read(4, "root:x:0:0::/root:/usr/bin/bash\n"..., 4096) = 1615
close(4)                                = 0
openat(AT_FDCWD, "/etc/group", O_RDONLY|O_CLOEXEC) = 4
read(4, "root:x:0:root\nsys:x:3:bin\nmem:x:"..., 4096) = 930
close(4)                                = 0
...
close(3)                                = 0
openat(AT_FDCWD, "/etc/localtime", O_RDONLY|O_CLOEXEC) = 3
read(3, "TZif2\0\0\0\0\0\0\0\0\0\0\0\0\0\0\0\0\0\0\17\0\0\0\17\0\0\0\0"..., 4096) = 2120
read(3, "TZif2\0\0\0\0\0\0\0\0\0\0\0\0\0\0\0\0\0\0\17\0\0\0\17\0\0\0\0"..., 4096) = 1317
close(3)                                = 0
write(1, "total 96\n-rw-r--r-- 1 sunflower "..., 444) = 444
close(1)                                = 0
close(2)                                = 0
+++ exited with 0 +++
\end{lstlisting}

From the trace, \texttt{ls -l} proceeds through four clear stages:
\begin{enumerate}
    \item \textbf{Dynamic Linking}: \texttt{ld-linux} opens \texttt{/etc/ld.so.cache}, \texttt{libcap.so.2}, and \texttt{libc.so.6} to map shared libraries into memory.
    \item \textbf{Locale and Directory Opening}: After probing \texttt{coreutils.mo} localization files, \texttt{ls} opens directory \texttt{.} with \texttt{O\_DIRECTORY} (receiving descriptor \texttt{3}).
    \item \textbf{Metadata and Identity Resolution}: To display owner names, group names, and local timestamps, \texttt{ls} consults \texttt{/etc/nsswitch.conf}, \texttt{/etc/passwd}, \texttt{/etc/group}, and \texttt{/etc/localtime}.
    \item \textbf{Output and Cleanup}: \texttt{ls} emits its formatted listing via \texttt{write(1, ...)} and closes standard streams \texttt{1} and \texttt{2} before exiting.
\end{enumerate}

\subsection{System Call Summary Profiling (\texttt{strace -c ls -l})}
We ran \texttt{strace -c ls -l >/dev/null} to obtain the aggregate count of all system calls invoked during one execution:
\begin{lstlisting}[style=consolestyle]
$ strace -c ls -l >/dev/null
% time     seconds  usecs/call     calls    errors syscall
------ ----------- ----------- --------- --------- ----------------
 31.16    0.000592           8        68        29 openat
 20.32    0.000386         386         1           execve
 12.68    0.000241          11        21           mmap
  8.53    0.000162           3        41           close
  7.11    0.000135           4        33           statx
  4.11    0.000078          13         6           mprotect
  2.47    0.000047           2        16           fstat
  2.32    0.000044           6         7           llistxattr
  2.05    0.000039           3        13           read
  1.79    0.000034           3         9           newfstatat
  1.42    0.000027           4         6         4 prctl
  0.95    0.000018           2         8           lseek
  0.79    0.000015           5         3           brk
  0.79    0.000015           3         4           rt_sigprocmask
  0.47    0.000009           4         2         2 ioctl
  0.47    0.000009           4         2           getdents64
  0.47    0.000009           4         2           getrandom
  0.42    0.000008           8         1         1 access
  0.37    0.000007           3         2           futex
  0.26    0.000005           5         1           write
  0.26    0.000005           5         1           prlimit64
  0.21    0.000004           4         1           arch_prctl
  0.21    0.000004           4         1           set_tid_address
  0.21    0.000004           4         1           set_robust_list
  0.16    0.000003           3         1           rseq
------ ----------- ----------- --------- --------- ----------------
100.00    0.001900           7       251        36 total
\end{lstlisting}

In the initial host desktop session (where the \texttt{systemd-userdbd} socket was reachable), \texttt{strace -c ls -l >/dev/null} recorded \textbf{260} total system calls (\textbf{30} \texttt{openat}, \textbf{28} \texttt{close}, \textbf{15} \texttt{read}, \textbf{1} \texttt{write}, plus socket/epoll IPC calls). In the isolated verification environment where \texttt{/run/systemd/userdb/} is unreachable, \texttt{libnss\_systemd} falls back to scanning filesystem \texttt{userdb} paths, yielding \textbf{251} total system calls (\textbf{123} filtered \texttt{openat}/\texttt{read}/\texttt{write}/\texttt{close} calls). Thus, one simple \texttt{ls -l} consistently executes roughly \textbf{250--260} system calls.

\subsection{Answers to Questions (Observe -- Part 1)}
\begin{enumerate}
    \item \textbf{Which system call does \texttt{ls} use to print to the screen? To which descriptor does it write?}
    \begin{itemize}
        \item \textbf{System call}: \texttt{write} (specifically \texttt{write(1, ...)}).
        \item \textbf{File descriptor}: \textbf{\texttt{1}} (\texttt{STDOUT\_FILENO}, standard output). On an interactive terminal, \texttt{ls} issues one \texttt{write(1, ...)} call per output line; when redirected or piped, \texttt{libc} block-buffers \texttt{stdout} and emits the full listing in one \texttt{write(1, ...)} call.
    \end{itemize}
\end{enumerate}

% ============================================================================
% 3. PART 2: A PROGRAM THAT CALLS THE KERNEL DIRECTLY
% ============================================================================
\section{Part 2: Calling the Kernel Directly (\texttt{syscalls.c})}

\subsection{Implementation Details}
We implemented \texttt{syscalls.c} using the POSIX system calls \texttt{open()}, \texttt{write()}, \texttt{read()}, and \texttt{close()} directly to create \texttt{note.txt}, write \texttt{"hello, kernel\textbackslash n"}, reopen the file read-only, and read the stored bytes back into user memory:

\begin{lstlisting}[style=cstyle, caption={Implementation of syscalls.c using POSIX system calls directly}]
/* syscalls.c - open, write, read and close, using only system calls.
 * Build:  gcc -Wall -Wextra -o syscalls syscalls.c
 * Run:    ./syscalls
 * See the calls: strace ./syscalls
 */
#include <fcntl.h>   /* open */
#include <unistd.h>  /* write, read, close */
#include <string.h>  /* strlen */
#include <stdio.h>   /* printf */

int main(void) {
    const char *msg = "hello, kernel\n";

    /* open() is a system call: ask the kernel for a file, get back a descriptor */
    int fd = open("note.txt", O_CREAT | O_WRONLY | O_TRUNC, 0644);
    if (fd < 0) { perror("open"); return 1; }
    printf("open returned descriptor %d\n", fd);

    /* write() is a system call: hand the bytes to the kernel */
    ssize_t n = write(fd, msg, strlen(msg));
    if (n < 0) { perror("write"); close(fd); return 1; }
    printf("write sent %zd bytes\n", n);
    close(fd);

    /* read it back, through the kernel again */
    fd = open("note.txt", O_RDONLY);
    if (fd < 0) { perror("open"); return 1; }
    char buf[64];
    n = read(fd, buf, sizeof(buf) - 1);
    if (n < 0) { perror("read"); close(fd); return 1; }
    buf[n] = '\0';
    printf("read got %zd bytes: %s", n, buf);
    close(fd);

    return 0;
}
\end{lstlisting}

\subsection{Compilation and Execution Output}
Compiling with \texttt{-Wall -Wextra} succeeded with zero warnings, and executing \texttt{./syscalls} produced the expected output and a 14-byte \texttt{note.txt} file:
\begin{lstlisting}[style=consolestyle]
$ gcc -Wall -Wextra -o syscalls syscalls.c
$ ./syscalls
open returned descriptor 3
write sent 14 bytes
read got 14 bytes: hello, kernel
$ wc -c note.txt
14 note.txt
\end{lstlisting}

\subsection{Tracing \texttt{./syscalls} with \texttt{strace}}
Tracing \texttt{openat}, \texttt{write}, \texttt{read}, and \texttt{close} during execution reveals both the dynamic loader setup and the exact six file-I/O calls from \texttt{main()}:
\begin{lstlisting}[style=consolestyle]
$ strace -e trace=openat,write,read,close ./syscalls
openat(AT_FDCWD, "/etc/ld.so.cache", O_RDONLY|O_CLOEXEC) = 3
close(3)                                = 0
openat(AT_FDCWD, "/usr/lib/libc.so.6", O_RDONLY|O_CLOEXEC) = 3
read(3, "\177ELF\2\1\1\3\0\0\0\0\0\0\0\0\3\0>\0\1\0\0\0\300y\2\0\0\0\0\0"..., 1024) = 1024
close(3)                                = 0
openat(AT_FDCWD, "note.txt", O_WRONLY|O_CREAT|O_TRUNC, 0644) = 3
write(1, "open returned descriptor 3\n", 27) = 27
write(3, "hello, kernel\n", 14)         = 14
write(1, "write sent 14 bytes\n", 20)   = 20
close(3)                                = 0
openat(AT_FDCWD, "note.txt", O_RDONLY)  = 3
read(3, "hello, kernel\n", 63)          = 14
write(1, "read got 14 bytes: hello, kernel"..., 33) = 33
close(3)                                = 0
+++ exited with 0 +++
\end{lstlisting}

Table~\ref{tab:syscall_map} summarizes the direct mapping between each C call in \texttt{syscalls.c} and the observed Linux kernel system call.

\begin{table}[H]
\centering
\small
\begin{tabularx}{\textwidth}{>{\raggedright\arraybackslash}p{5.5cm}>{\raggedright\arraybackslash}Xc}
\toprule
\textbf{C Source Statement} & \textbf{Observed Kernel System Call} & \textbf{Return} \\
\midrule
\texttt{open("note.txt", O\_CREAT | O\_WRONLY | O\_TRUNC, 0644)} & \texttt{openat(AT\_FDCWD, "note.txt", O\_WRONLY|O\_CREAT|O\_TRUNC, 0644)} & \texttt{3} \\
\texttt{write(fd, msg, strlen(msg))} & \texttt{write(3, "hello, kernel\textbackslash n", 14)} & \texttt{14} \\
\texttt{close(fd)} & \texttt{close(3)} & \texttt{0} \\
\texttt{open("note.txt", O\_RDONLY)} & \texttt{openat(AT\_FDCWD, "note.txt", O\_RDONLY)} & \texttt{3} \\
\texttt{read(fd, buf, sizeof(buf) - 1)} & \texttt{read(3, "hello, kernel\textbackslash n", 63)} & \texttt{14} \\
\texttt{close(fd)} & \texttt{close(3)} & \texttt{0} \\
\bottomrule
\end{tabularx}
\caption{Mapping between C library calls in \texttt{syscalls.c} and kernel system calls observed via \texttt{strace}.}
\label{tab:syscall_map}
\end{table}

Running \texttt{strace -c ./syscalls >/dev/null} confirms that the entire program executes only \textbf{41} total system calls across 18 distinct syscall types from startup to exit.

\subsection{Answers to Questions (Observe -- Part 2)}
\begin{enumerate}
    \item \textbf{The descriptor returned by the first \texttt{openat} is 3, not 0, 1 or 2. Why are those three already taken?} \\
    When any Linux/UNIX process is created and started by a shell, three standard file descriptors are already open by convention:
    \begin{itemize}
        \item \textbf{\texttt{0} (\texttt{STDIN\_FILENO})}: Standard input (\texttt{stdin}), attached to the terminal keyboard or input stream.
        \item \textbf{\texttt{1} (\texttt{STDOUT\_FILENO})}: Standard output (\texttt{stdout}), attached to the terminal screen or redirected output target.
        \item \textbf{\texttt{2} (\texttt{STDERR\_FILENO})}: Standard error (\texttt{stderr}), attached to the terminal screen for diagnostics.
    \end{itemize}
    Because POSIX mandates that \texttt{open()}/\texttt{openat()} return the \textbf{lowest-numbered unused descriptor}, the first newly opened file (\texttt{note.txt}) receives descriptor \textbf{\texttt{3}}. After \texttt{close(3)} frees slot \texttt{3}, the second \texttt{open("note.txt", O\_RDONLY)} immediately reuses descriptor \textbf{\texttt{3}}.
\end{enumerate}

% ============================================================================
% 4. PART 3: COUNT ONE KIND OF CALL
% ============================================================================
\section{Part 3: Counting One Kind of System Call}

\subsection{Filtering and Counting Specific Syscalls}
We isolated and counted individual system calls (\texttt{write} and \texttt{openat}) for both \texttt{ls -l} and \texttt{./syscalls}:
\begin{lstlisting}[style=consolestyle]
$ strace -f -e trace=write ls -l 2>&1 >/dev/null | grep -c '^write'
1

$ strace -f -e trace=openat ls -l 2>&1 >/dev/null | grep -c '^openat'
68

$ strace -e trace=write ./syscalls 2>&1 >/dev/null | grep -c '^write'
2
\end{lstlisting}

\subsection{Analysis of Buffering and Output Redirection}
The command pipeline \texttt{2>\&1 >/dev/null | grep -c '\^{}write'} performs two subtle operations:
\begin{enumerate}
    \item \textbf{Separating Program Output from \texttt{strace} Output}: Because \texttt{strace} prints trace lines to \texttt{stderr} (fd \texttt{2}) while the traced program prints to \texttt{stdout} (fd \texttt{1}), \texttt{2>\&1} sends \texttt{strace}'s \texttt{stderr} into the pipe while \texttt{>/dev/null} discards the program's \texttt{stdout}. Without \texttt{>/dev/null}, \texttt{./syscalls}'s own output line \texttt{"write sent 14 bytes"} matches \texttt{\^{}write} in \texttt{grep}, falsely inflating the count from \texttt{2} to \texttt{3}.
    \item \textbf{TTY Line Buffering vs.\ Redirected Block Buffering}: Redirecting \texttt{stdout} to \texttt{/dev/null} switches \texttt{libc}'s \texttt{stdout} stream from line buffering to block buffering:
    \begin{itemize}
        \item In \texttt{./syscalls} with \texttt{stdout} redirected to \texttt{/dev/null}, all three \texttt{printf()} messages accumulate in the user-space buffer and flush in \textbf{1} \texttt{write(1, ..., 80)} call on exit, plus \textbf{1} direct \texttt{write(3, "hello, kernel\textbackslash n", 14)} call, totaling \textbf{2} \texttt{write} syscalls. On a live terminal (TTY), \texttt{stdout} flushes after every \texttt{\textbackslash n}, producing \textbf{4} \texttt{write} syscalls (3 \texttt{write(1, ...)} + 1 \texttt{write(3, ...)}).
        \item Similarly, \texttt{ls -l} issues \textbf{1} \texttt{write(1, ...)} call when redirected to \texttt{/dev/null}, versus one \texttt{write(1, ...)} per line on a terminal.
    \end{itemize}
\end{enumerate}

\begin{table}[H]
\centering
\small
\begin{tabularx}{\textwidth}{>{\raggedright\arraybackslash}p{1.9cm}>{\raggedright\arraybackslash}p{1.3cm}>{\centering\arraybackslash}p{2.3cm}>{\centering\arraybackslash}p{2.3cm}>{\raggedright\arraybackslash}X}
\toprule
\textbf{Command} & \textbf{Syscall} & \textbf{Redirected (\texttt{>/dev/null})} & \textbf{TTY (\texttt{strace -o})} & \textbf{Explanation} \\
\midrule
\texttt{ls -l} & \texttt{write} & 1 & 8 (header + 7 entries) & Block-buffered flush vs.\ per-line TTY output \\
\texttt{ls -l} & \texttt{openat} & 68 / 30 (host) & 68 / 30 (host) & Shared libs, locales, \texttt{.}, NSS/userdb, \texttt{/etc/localtime} \\
\texttt{./syscalls} & \texttt{write} & 2 & 4 & 1 merged \texttt{printf} flush vs.\ 3 TTY flushes, plus 1 \texttt{write(3, ...)} \\
\bottomrule
\end{tabularx}
\caption{Comparison of measured system call counts under output redirection and interactive terminal execution.}
\label{tab:part3_counts}
\end{table}

% ============================================================================
% 5. CONCLUSIONS
% ============================================================================
\section{Conclusions}

\begin{enumerate}
    \item \textbf{Visibility across the System Call Boundary}: \texttt{strace} exposes every transition between user space and the Linux kernel, revealing that a routine \texttt{ls -l} invocation executes roughly 250--260 system calls spanning dynamic linking, localization, directory traversal, identity resolution, and terminal output.
    \item \textbf{Direct POSIX File I/O}: Writing \texttt{syscalls.c} with \texttt{open()}, \texttt{write()}, \texttt{read()}, and \texttt{close()} confirmed how user-space C calls map directly to kernel syscalls (\texttt{openat}, \texttt{write}, \texttt{read}, \texttt{close}) without buffered \texttt{FILE*} wrappers.
    \item \textbf{Reserved Descriptors and Reuse}: Descriptors \texttt{0}, \texttt{1}, and \texttt{2} are permanently reserved at startup for \texttt{stdin}, \texttt{stdout}, and \texttt{stderr}; consequently, the kernel assigns descriptor \texttt{3} to the first opened file and recycles slot \texttt{3} once closed.
    \item \textbf{Buffering Efficiency}: Counting \texttt{write} calls under redirection versus interactive TTY execution demonstrated how \texttt{libc} stream buffering merges multiple \texttt{printf()} outputs into a single kernel transition, reducing system call overhead.
\end{enumerate}

\end{document}
